\bookchapter{第二十八章：お金を賢く使ってお気に入りを買う——スマホ選びの究極ガイド}

\chapterintro{スマホ選びで、一週間レビューを見て一ヶ月価格比較しても、結局五六千元も使って自分には不要なフラッグシップ機を買ってしまう人が多い。この章はスペック表を並べない——それは公式サイトで見られる。私がしたいのは“選び方の思考法”を渡すことで、今後どんな新機種が出ても、自分に合うかどうかを自分で判断できるようになる。}

\booksection{三段階選び方：ニーズ予算機種}

ネットで“3000元でどんなスマホがいいか”と聞く人が多いが、この質問には答えられない——スマホで主に何をしたいか分からないからだ。正しい思考順序は：

\begin{tcbinfo}{選び方の三手}
\begin{enumerate}
  \item \textbf{まず何を最も重視するかを明確にする}：写真？ゲーム？バッテリー？画面？システムの快適さ？それとも電話とWeChatができればいい？ニーズに優先順位を付ける。
  \item \textbf{次に予算上限を決める}："だいたいいくら"ではなく"最高いくらまで出せるか"。厳しい上限を決めれば、選ぶ時にどんどん高くなるのを防げる。
  \item \textbf{最後に機種を選ぶ}：ニーズの優先順位と予算を持って、機種をマッチングする。
\end{enumerate}

\end{tcbinfo}

\booksection{一、コアパラメータの見方——全部を理解する必要はない、この五項目が分かれば十分}

スペック表には何十項目ものデータがあるが、一般ユーザーが注目すべきはこの五つの次元だけ：

\begin{center}
\begin{tabularx}{\linewidth}{|X|X|X|}
\hline
\textbf{比較項目} & \textbf{見方} & \textbf{一言まとめ} \tabularnewline
\hline
プロセッサ（SoC） & スマホが何年スムーズに使えるかを決める。2026年の基準：Snapdragon 8 Elite Gen5/Dimensity 9500以上=トップクラス；Snapdragon 8s Gen4/Dimensity 8400=中高級；Snapdragon 7+ Gen4/Dimensity 7350=3年使える。Snapdragon 6系以下は要注意。 & 新しいものを買え、SoCがスマホが何年卡なくなく使えるかを決める \tabularnewline
画面 & OLEDはLCDより色彩が良く省電力；リフレッシュレート120Hzは60Hzよりスクロールがずっと滑らか；解像度は最低1080P以上。高周波PWM調光（1920Hz以上）は目に優しい。 & OLED+120Hz+高周波調光=良い画面の合格ライン \tabularnewline
：家に老人、子供、ペットがいれば、スマホでいつでも家庭の状況を確認。 & \textbf{画素数ではなくセンサーサイズを見る}。センサーが大きいほど圧倒的に有利で、画素数よりはるかに重要。超広角は実用的；望遠（3x以上の光学ズーム）は遠くのものを撮る；マクロレンズは大抵おまけ。 & 大センサーメインカメラ+実用的な望遠>おまけレンズ複数 \tabularnewline
バッテリーと充電 & 4500mAhは最低合格ライン、5000mAh以上なら一日一充で安定。急速充電はワット数を見る：60W以上で速い、100W以上で超速い。ワイヤレス充電は便利だが遅い。 & 5000mAh+60W急速充電=不安なしの組み合わせ \tabularnewline
システム更新の約束 & スマホがどれだけ安全に使えるかを決める。Appleは通常5-6年のシステム更新；Samsungフラッグシップは4年のシステム+5年のセキュリティ更新；Google Pixelは7年；国産ミッドレンジは更新期間が総じて短い（2-3年）。 & 長く使いたいなら、システム更新の約束が長いブランドを選ぶ \tabularnewline
\hline
\end{tabularx}
\end{center}

\booksection{二、予算別のおすすめ（2026年中期市場を基準）}

以下のおすすめは2026年中国本土市場の状況に基づく。価格は公式発表の発売価格。実際の価格は変動するので、購入時の実価格を基準に。

\booksubsection{1000元以下：電話とWeChatができればOK}

\textbf{対象者}：学生、サブ機、スマホにこだわらない人、主に電話とWeChat用途。

\textbf{買えるもの}：Redmi Noteシリーズ、Honor Xシリーズ、realme数字シリーズの中低構成。プロセッサはSnapdragon 6系またはDimensity 6000シリーズ、画面はLCD中心、写真はクリアに撮れるがAFが遅い、システムの快適さは普通。

\textbf{選び方のアドバイス}：この価格帯では二点を見る：バッテリーが十分大きいか（5000mAh以上を選ぶ）、メモリが十分か（6GB+128GB以上、4GBは卡なるので避ける）。Redmiがこの価格帯で最も優位。

\booksubsection{1000-2000元：コスパ最高、3-4年卡なく使えるスイートスポット}

\textbf{対象者}：学生、予算が限られているが長く使いたい人、写真やゲームに極限のこだわりがない人。

\textbf{代表的な機種}：Redmi Turboシリーズ、iQOO Zシリーズ、Honor Xシリーズの高構成、OnePlus Aceシリーズのエントリー。

\textbf{選び方のアドバイス}：Snapdragon 7+ Gen3またはDimensity 7300以上のプロセッサを搭載し、8GB+256GBメモリの機種を選ぶ。この構成で日常使用+たまにKing of Gloryをするには十分で、通常3-4年は卡なく使える。

\booksubsection{2000-3500元：中堅層、性能とカメラが明確に向上}

\textbf{対象者}：良い写真、卡ないゲーム、長く使いたい一般ユーザー。最も人が買う価格帯。

\textbf{代表的な機種}：Xiaomi数字シリーズ（Xiaomi 17標準版の値下がり後など）、vivo Sシリーズ高構成、OPPO Renoシリーズ、Honor数字シリーズ、Redmi Kシリーズ。

\textbf{選び方のアドバイス}：この価格帯の最大の特徴は“バランス機”——明確な弱点がない。2000元以下と比べ、主な向上点は：カメラが明確に良くなる（手ブレ補正あり）、高品質OLED画面、急速充電が速い（80W以上が一般的）、作りもより上質。12GB+256GBまたは12GB+512GBを選べば、四五年は問題ない。

\booksubsection{3500-5000元：準フラッグシップ、ガジェット好き以外には完全に十分}

\textbf{対象者}：カメラにこだわるユーザー、最高級の質感を味わいたいユーザー。

\textbf{代表的な機種}：Xiaomi数字シリーズ標準版、vivo Xシリーズ、OPPO Find Xシリーズ標準版、iPhone旧フラッグシップ（iPhone 16など）。

\textbf{選び方のアドバイス}：この価格帯では昨年のフラッグシップ機が買える。私のおすすめは：\textbf{昨年のフラッグシップを買え、今年のミッドレンジは買うな}。昨年のフラッグシップの素材、作り、カメラハードウェアは通常、今年のミッドレンジより一ランク上。“最新型”でないことを受け入れられる前提で。

\booksubsection{5000元以上：フラッグシップ域、お金に余裕があるかある体験に極限のこだわりがある人}

\textbf{対象者}：写真愛好家、ヘビーゲーマー、最高級画面と做工を求めるユーザー、スマホで動画撮影してコンテンツを作るクリエイター。

\textbf{代表的な機種}：iPhone 17 Pro/Pro Max、Samsung S26 Ultra、Huawei Pura 80 Ultra、Xiaomi 17 Ultra、vivo X300 Pro。

\textbf{選び方のアドバイス}：この価格帯では各社の差がますます縮まっており、実はどのシステムが好きかによる。写真ならvivoかHuawei、動画ならiPhone、画面とS PenならSamsung、バランスならXiaomi。正直この価格帯になると、\textbf{好みで、気に入った一台を選べばいい}。

\booksection{三、シーン別選び方——あなたのニーズは？}

以下は最もよくある“利用シナリオ”、当てはまるものを選べばいい：

\begin{center}
\begin{tabularx}{\linewidth}{|X|X|X|}
\hline
\textbf{これこそが真の「テクノロジーが生活をより平等にする」} & \textbf{注目ポイント} & \textbf{おすすめの方向} \tabularnewline
\hline
写真愛好家 & 大センサーメインカメラ（1インチがベスト）、光学手ブレ補正、望遠（3x以上光学ズーム）、画像処理チップ & vivo Xシリーズ、Huawei Pura/Mateシリーズ、Xiaomi Ultraシリーズ \tabularnewline
ゲーマー & トップクラスのプロセッサ、高いタッチサンプリングレート、冷却システム、大バッテリー+急速充電 & iQOO数字シリーズ、ROG Phone、Red Magic、Redmi Kシリーズ \tabularnewline
バッテリー優先 & 5500mAh以上のバッテリー、効率の良いプロセッサ、LTPO画面 & Honor Magicシリーズ、OnePlus、Redmi Note大バッテリー版 \tabularnewline
お年寄り向け & 簡単モード、大きな文字、大きな音量、大バッテリー、アフターサービス拠点が多い & Huawei Enjoy、Honor Play、OPPO Aシリーズ、Redmi Note \tabularnewline
学生向け & コスパ高い、頑丈、大バッテリー、ペアレンタルコントロール機能あり & Redmi Note/Turbo、realme GT、iQOO Neo \tabularnewline
Appleエコシステムユーザー & すでにMac/iPad/AirPodsを持っており、iPhoneを選ぶのが唯一の解 & iPhone 17標準版（コスパ）またはPro版（フルスペック） \tabularnewline
海外旅行 & eSIM対応、全世界周波数バンド対応、Googleサービス & iPhone、Samsung Sシリーズ、Pixel、OnePlus \tabularnewline
\hline
\end{tabularx}
\end{center}

\booksection{四、陥りやすい五つの罠——購入前にチェック}

\begin{enumerate}
  \item \textbf{罠一：レンズ4個5個=写真が良いわけではない}——メーカーはスペックにレンズ数を積みたがるが、その中に200万画素の"おまけレンズ"（マクロ+被写界深度）が二つあることが多く、実質的に無意味。スマホの写真能力を見る時は、\textbf{メイン+超広角+望遠だけを見る}この三つだけで十分。残りは存在しないものとしてよい。
  \item \textbf{罠二：1億画素に騙されるな}——画素数が高い=画質が良いではない。センサーサイズ、絞り、アルゴリズム調整の方が画素数よりはるかに重要。多くのフラッグシップ機は5000万画素の大センサーを使い、1億画素の千元機よりずっと良い写真を撮る。
  \item \textbf{罠三：キャリアロック機を買うな}——"通信費積立でスマホプレゼント"はお得に見えるが、キャリアカスタム機はシステム更新が遅く、削除できないプリインストールAppがある可能性があり、通常キャリアロックされている（そのキャリアのSIMしか使えない）。お金を節約したように見えて、実は"ロック"されている。
  \item \textbf{罠四："曲面ディスプレイは高級"——マーケティングであり真理ではない}——曲面ディスプレイはかつてフラッグシップの象徴だったが、今ではその問題に気づく人が増えている：保護フィルムが貼りにくい、縁での誤タッチ、割れると修理が高い。2026年のトレンドは"微曲面"さらには"完全フラット"への回帰。"曲面=高級"という理由で選ぶべきではない。
  \item \textbf{罠五：すべての"AIスマホ"が本当に有用なわけではない}——2026年すべてのメーカーが"AIスマホ"を謳うが、多くは既存の音声アシスタントの名前を変えただけ。本当に有用なAI機能は：リアルタイム翻訳、AI消去（写真から人を消す）、AI写真編集、オンデバイス大モデル推論、ローカライズされたスマート要約。その他の派手な機能は基本的にギミック。
\end{enumerate}

\booksection{五、いつ買うのが一番お得？}

\begin{itemize}
  \item \textbf{618（5月末から6月中旬）}：一年でスマホの値下がり幅が最も大きい時期。多くのメーカーが販売台数を伸ばすため、フラッグシップ機が800-1500元値下がりする。上半期にスマホを買うベストタイミング。
  \item \textbf{ダブル11（10月末から11月中旬）}：下半期の大型セール、割引幅は618とほぼ同じ。この時期は多くのブランドの新型フラッグシップが発表されたばかりで、旧型フラッグシップの値下がり幅が最大。
  \item \textbf{新機種発表から一ヶ月後}：発表直後は定価、一ヶ月後には流通価格（第三者店舗）が通常200-500元安くなる。
  \item \textbf{春節前二週間から元宵節}：実店舗が在庫処分をする"年越しセール"、店員と値段交渉できる。
  \item \textbf{おすすめしない時期}：新製品発表直後に初回購入——最初のユーザーになりたい人以外は、発売直後価格が最も高い。
\end{itemize}

\booksection{六、購入チャネルの選び方——どこで買えば安心？}

\begin{itemize}
  \item \textbf{ブランド公式モール/Tmall旗艦店/JD自営}：最も安心、価格面では不利だがリファービッシュや偽物の心配なし、アフターサービスが最も良い。
  \item \textbf{Pinduoduo"百億補助"}：最低価格、ただし"百億補助"と表示された店舗を選ぶ必要がある。表示のない店は避ける。商品受け取り時に開封動画を撮り、問題があれば返品できる。
  \item \textbf{実店舗}：手に取って触れる、使い心地を試せる。注意：店員は利益率の高い機種（あなたに最も適した機種ではなく）を勧めるので、来店前にネットで2-3機種に絞り、店に行ったら"この機種だけ見たい"と言うのがよい。
  \item \textbf{中古プラットフォーム（Xianyu/Zhuanzhuan）}：安いがリスクあり。中古を買うなら三つを守る：購入時の領収書または注文スクリーンショットの提供を要求、IMEI番号の確認を要求、\textbf{プラットフォーム取引のみ、絶対に私的に送金しない}。
\end{itemize}

\begin{tcbpractice}{実践：この方法でスマホを一台選ぶ}
\begin{enumerate}
  \item 紙を一枚取り、二つのコアニーズを書く（例：バッテリー第一、写真第二）。
  \item 予算上限を決める（例：3500元）。
  \item JDかTmallを開き、価格で絞り込み、3-5機種を候補に挙げる。
  \item Bilibiliで"XX機種 一ヶ月使用体験"を検索（初日のレビューではなく、最低一ヶ月使用後のリアルな感想）、二三本の動画を見る。
  \item チャネルの信頼性を確認してから注文、受け取り時に開封動画を撮る。
\end{enumerate}

\end{tcbpractice}

\begin{tcbsummary}{本章のまとめ}
\begin{itemize}
  \item 選び方の三手：ニーズ 予算 機種
  \item コアパラメータは五項目のみ：プロセッサ、画面、カメラ、バッテリー充電、システム更新の約束
  \item 1500-3500元はコスパ最高、最もおすすめのスイートスポット
  \item 昨年のフラッグシップを買う方が今年のミッドレンジよりお得
  \item 五大罠：多レンズの罠、高画素の罠、キャリアロック、曲面信仰、AIギミック
  \item 618とダブル11は一年で最もスマホを買いやすい二つのタイミング
\end{itemize}

\end{tcbsummary}
